\paragraph{Generalization requirements.} This evaluation covers one QSP model in one disease area, extracted and curated by one modeling group, so it characterizes the framework rather than measuring how broadly it generalizes. Several components are domain-independent: the separation of an inputs layer from a calibration layer, the typed-role fields, the value-in-snippet, DOI-resolution, and code-execution validators, and the provenance requirements all apply to any literature-derived parameter. The forward-model type library is the main domain-specific surface. \mtAlgebraicPct{} of SubmodelTargets used the generic algebraic type, which indicates that the 15 built-in templates covered roughly half of the use cases even within this model; a portion of those algebraic uses are genuinely simple derived relations (half-lives, ratios) for which no structured template is needed, but the fraction also shows the template library is not exhaustive. A non-oncology application, for instance a PK/PD model, would likely require additional templates for compartmental pharmacokinetics, receptor occupancy, or indirect-response dynamics; the generic algebraic, direct-fit, and custom fallbacks accept such cases today but without the structured role typing that the built-in templates provide. The model-aware search and shared-parameter-name binding assume a structured QSP model is already in hand, which holds in any domain but means the framework supports calibration of an existing model rather than model construction.
